\subsection{Asymptotic behavior at large scale when $n\ge4$}\label{large-scale_asympt} 
%
%
In this subsection we investigate the asymptotic behavior of the magnitude function $M_X(t)$ at large scale ($t\to+\infty$) as preparation for the proof when $n\ge4$. For this purpose it seems that the use of the formal magnitude $m_X(q)$ $(q=e^{-t})$ would make the description easier to read. 
Assume $n\ge4$ in what follows. 

(i) First we show that a finite rationally independent metric space is determined by the triples of lengths of three consecutive edges that form triangles or open $3$-paths. 


Let $X=\{P_1,\dots,P_n\}$. 
Let $\widetilde{\mathcal{S}}_\triangle$ (or $\widetilde{\mathcal{S}}_\otp$) be the multiset of the multisets of lengths of edges forming triangles (or respectively, open $3$-paths): 
\begin{equation}
\begin{array}{rcl}
\widetilde{\mathcal{S}}_\triangle&=&\displaystyle \left[\,[d_{ij},d_{jk},d_{ki}]\,|\,(i,j,k)\in \mathcal{I}_\triangle \right], \notag \\[1mm]
%
\widetilde{\mathcal{S}}_\otp&=&\displaystyle \left[\,[d_{ij},d_{jk},d_{kl}]\,|\,(i,j,k,l)\in \mathcal{I}_\otp \right].
\notag 
\end{array} 
\end{equation}
%
Note that $\widetilde{\mathcal{S}}_\triangle$ and $\widetilde{\mathcal{S}}_\otp$ have no information about the subscripts of $d_{\ast\ast}$; in other words, even if we know the triplet of edge lengths, we do not know their vertices.
%
\begin{lemma}\label{lemma_determine}
A finite $3$-generic metric space is determined by $\widetilde{\mathcal{S}}_\triangle$ and $\widetilde{\mathcal{S}}_\otp$. 
\end{lemma}

\begin{proof}
The $3$-genericity condition ($\mbox{g}_3$) implies that the edge lengths are different from each other. Therefore from $\widetilde{\mathcal{S}}_\triangle$ %
we can obtain the multiset of the edge lengths $\widetilde{\mathcal{S}}_1=[d_{ij}\,|\,1\le i<j\le n]$ (in fact it is a set in this case). 

Let $\widetilde{\mathcal{S}}_\sotp$ (or $\mathcal{S}_\sotp$) be the multiset of the multisets of lengths of edges forming simple open $3$-paths (or respectively, with the information of the length of the middle edge):
\begin{eqnarray}
%
%
\widetilde{\mathcal{S}}_\sotp&=&\displaystyle \left[\,[d_{ij},d_{jk},d_{kl}]\,|\,(i,j,k,l)\in \mathcal{I}_\sotp \right],
\notag \\[1mm]
%
\mathcal{S}_\sotp&=&\displaystyle \left[(d_{jk}, [d_{ij},d_{kl}])\,|\,(i,j,k,l)\in \mathcal{I}_\sotp \right]. \notag
\end{eqnarray}

First remark that $\widetilde{\mathcal{S}}_\sotp$ can be obtained from $\widetilde{\mathcal{S}}_\otp$ by removing multisets with duplications. 

Next remark that the data of $\widetilde{\mathcal{S}}_\triangle$ and $\widetilde{\mathcal{S}}_\sotp$ produce $\mathcal{S}_\sotp$, namely, the information of the middle edges of simple open $3$-paths can be obtained from $\widetilde{\mathcal{S}}_\triangle$ and $\widetilde{\mathcal{S}}_\sotp$. 
This is because the two edges at the ends of a simple open $3$-path cannot form a triangle with another edge. 

Finally we show that a finite $3$-generic metric space is determined by $\widetilde{\mathcal{S}}_1, \widetilde{\mathcal{S}}_\triangle$ and $\mathcal{S}_\sotp$. Suppose we have the data of $\widetilde{\mathcal{S}}_1, \widetilde{\mathcal{S}}_\triangle$ and $\mathcal{S}_\sotp$. Choose a triangle and an edge of it. 
We may label the three vertices $P_1, P_2$ and $P_3$ so that the edge we selected is $P_1P_2$. Let $\ell_\a=d(P_2,P_3), \ell_\beta=d(P_3,P_1)$ and $\ell_\ga=d(P_1,P_2)$. 
There are still $n-3$ triangles containing the edge $P_1P_2$. 
There are two ways to attach each triangle to $P_1P_2$, but one is determined from the information of open 3-paths and middle edges as follows. 
Suppose $[\ell_\la,\ell_\mu,\ell_\ga]\in\widetilde{\mathcal{S}}_\triangle$. Let the remaining vertex be $P_4$, say. 
Note that both $\ell_\la, \ell_\mu, \ell_\a$ and $\ell_\la, \ell_\mu, \ell_\beta$ form open $3$-paths. 
If the middle edge of $\ell_\la, \ell_\mu, \ell_\a$ is $\ell_\la$ then the triangle $\triangle P_4P_1P_2$ is attached to the edge $P_1P_2$ in a way that $d(P_4,P_1)=\ell_\mu$ and $d(P_4,P_2)=\ell_\la$, and if not the other way. 

After attaching the remaining $n-4$ triangles to $P_1P_2$, label the remaining vertices $P_5,\dots,P_n$. The lengths $d(P_i,P_1), d(P_i,P_2)$ $(4\le i \le n)$ are determined by the procedure described above. 
%
The length $d(P_i,P_j)$ $(3\le i<j\le n)$ is determined as the unique element $\ell$ in $\widetilde{\mathcal{S}}_1$ such that $[\ell, d(P_i,P_1), d(P_j,P_1)]$ is an element of $\widetilde{\mathcal{S}}_\triangle$. 
%
\end{proof}


\medskip
(ii) Next we prepare a proposition which we will use to get a multiset consisting of the sums of lengths of edges forming triangles and a multiset consisting of the sums of lengths of edges forming open $3$-paths from the formal magnitude $m_X(q)$. 



\begin{proposition}\label{Leinster}{\rm (Leinster \cite{L19})} 
The formal magnitude of a finite metric space $X$ is given by 
\begin{equation}\label{eqL}
m_X(q)=\sum_{k=0}^\infty(-1)^k\sum_{(i_0,\dots,i_k)\in\mathcal{I}_k}q^{d_{i_0i_1}+\dots+d_{i_{k-1}i_k}}.
\end{equation}
\end{proposition}
%
In fact it was proved in Proposition 3.9 of \cite{L19} for graphs, and the proof given there works for any finite metric space as well (cf. formula (1) of \cite{HW} and Corollary 7.15 of \cite{LS}). 

%
\begin{definition}\label{def_d-index}  \rm 
\begin{enumerate}
\item 
Let $\mathcal{P}_{all}$ be the commutative monoid generated by $\widetilde{\mathcal{S}}_1=[d_{ij}]$;
\begin{equation}\label{P_all}
\mathcal{P}_{all}=\Big\{\sum_{i<j}a_{ij}d_{ij}\,|\,a_{ij}\in\bar\NN%
\Big\}.
%
\end{equation}
%
\item Let $\mathcal{P}$ $(\mathcal{P}\subset\mathcal{P}_{all})$ be the set of positive exponents that appear in $m_X(q)$. 
%
%
\item 
We define the {\em $d$-index} of a term $q^{\,\sum_{i<j}a_{ij}d_{ij}}$ to be $\sum_{i<j}a_{ij}$. 
%
\end{enumerate}
\end{definition}
%

%
Let $m_{X,3}(q)$ be the sum of all the terms in $m_X(q)$ with the $d$-index less than or equal to $3$; 
\[
m_{X,3}(q)=\sum_{k=0}^3(-1)^k\sum_{(i_0,\dots,i_k)\in\mathcal{I}_k}q^{d_{i_0i_1}+\dots+d_{i_{k-1}i_k}}.
\]
%
By dividing $2$- and $3$-paths into closed paths and open paths we obtain 
%
\begin{eqnarray}
m_{X,3}(q)&=&\displaystyle n-2\sum_{i<j}q^{d_{ij}}+2\sum_{i<j}q^{2d_{ij}}
+2\sum_{\mathcal{I}_\scotwop} q^{d_{ij}+d_{jk}}
%
%
%
%
-6\sum_{\mathcal{I}_\triangle} q^{d_{ij}+d_{jk}+d_{ki}} 
%
-2\sum_{\mathcal{I}_\scotp} q^{d_{ij}+d_{jk}+d_{kl}}.
%
%
%
%
%
\label{m_3}
\end{eqnarray}
%
%



\medskip
(iii) Finally we give a lemma to obtain the exponents and coefficients %
of $m_X(q)$ %
from the magnitude function $M_X(t)$. 

\begin{lemma}\label{lemma_P} 
Suppose $m_X(q)$ is expressed as 
\[
m_X(q)=\sum_{m=0}^\infty a_mq^{\a_m} \quad (a_m\in\RR,\,\{\a_0,\a_1.\dots\}=\mathcal{P},\,\a_0<\a_1<\dots ).
\]
Then $\a_0=0$ and $\displaystyle a_0=\lim_{t\to+\infty}M_X(t)=\#X$, and $\a_m$ and $a_m$ are given inductively by 
\[
\begin{array}{rcl}
\a_m&=&\displaystyle \lim_{t\to+\infty}\frac{\displaystyle \left|\log \left(M_X(t)-\sum_{i=0}^{m-1}a_i\, e^{-t\a_i}\right)\right|}{t}, \\[6mm]
%
a_m&=&\displaystyle \lim_{t\to+\infty} e^{t\a_m}\left(M_X(t)-\sum_{i=0}^{m-1}a_i\, e^{-t\a_i}\right).
\end{array}
\]

\end{lemma}


%
\subsection{Proof of Part (2) of Theorem \ref{theorem}}\label{part2} 
%
%
By Lemma \ref{lemma_P} we obtain the multiset $\mathcal{P}$ from $M_X(t)$. 
From $\mathcal{P}$ we can obtain the multiset of the edge lengths $\widetilde{\mathcal{P}}_1=[d_{12},\dots,d_{n-1n}]$ as follows. 
Remark that the rational independence condition (ri) implies that $d_{ij}$'s are all different from each other, hence $\widetilde{\mathcal{P}}_1$ is an ordinary set in fact. 
Define $\ell_1,\dots,\ell_N$ inductively, where $N={n\choose2}$, by 
%
\begin{equation}\label{how_to_determine_ell}
\begin{array}{l}
\displaystyle \ell_1=\min\mathcal{P}, \\[2mm]
\displaystyle \ell_2=\min\left(\mathcal{P}\setminus \{\tau\,\ell_1\,|\,\tau\in\NN\}\right), \\[2mm] 
\displaystyle \ell_3=\min\left(\mathcal{P}\setminus \{\tau_1\,\ell_1+\tau_2\,\ell_2\,|\,\tau_1,\tau_2\in\NN\cup\{0\}\}\right), \\[2mm]
\dots
\end{array}
\notag
\end{equation}
%
Then the (multi)set of edge lengths $\widetilde{\mathcal{P}}_1$ is given by the (multi)set $[\ell_1,\dots,\ell_N]$. 

Put 
\[\begin{array}{rcl}
\widetilde{\mathcal{S}}_p&=&\displaystyle [\,[\ell_{\a_{i_1}},\dots,\ell_{\a_{i_p}}]\,|\,1\le \a_{i_1}\le\dots\le\a_{i_p}\le N] \qquad (p\in\NN), \\[2mm]
\widetilde{\mathcal{S}}&=&\displaystyle \bigcup_{p\in\NN}\widetilde{\mathcal{S}}_p.
\end{array}
\]
Note that $\widetilde{\mathcal{S}}_1=\widetilde{\mathcal{P}}_1$ and that $\widetilde{\mathcal{S}}_3\supset \widetilde{\mathcal{S}}_\triangle, \widetilde{\mathcal{S}}_\otp$. 
The rational independence condition (ri) implies that a map $\Sigma\colon\widetilde{\mathcal{S}}\to\RR$ given by 
\[\Sigma\left(\,[\ell_{\a_{i_1}},\dots,\ell_{\a_{i_p}}]\,\right)=\ell_{\a_{i_1}}+\dots+\ell_{\a_{i_p}}\]
is injective. 
Since the ``triangle'' terms and the ``{open $3$-path}'' terms in \eqref{m_3} have different coefficients, by comparing $\mathcal{P}$ and $\Sigma(\widetilde{\mathcal{S}}_3)$ %
we obtain $\Sigma(\widetilde{\mathcal{S}}_\triangle)$ and $\Sigma(\widetilde{\mathcal{S}}_\otp)$.
Since $\Sigma$ is injective, the conclusion follows from Lemma \ref{lemma_determine}. 
